\section{比较膜量子化与几何量子化}\label{comparing}

\subsection{复余切丛与形变量子化}\label{cotangent}

几何量子化与膜量子化，是为量子化辛流形 $M$ 引入必要额外结构的两种不同方式。这两大量子化方法在什么情况下一致？我们能期望的最好答案，是找到 $M$ 的复化 $Y$ 与 $M$ 的极化 $\P$ 上的条件，使得基于 $Y$ 的膜量子化等价于基于 $\P$ 的几何量子化。本节将针对第 \ref{gq} 节引入的两类基本极化——$M$ 的实极化与复极化——解释这样的等价判据。

先描述人们可能期望比较膜量子化与几何量子化的最明显情形。设 $M=T^*N$ 是余切丛，带自然辛结构 $\omega=\sum_i\d p_i\d q^i$，其中 $q^i$ 为 $N$ 的局部坐标，$p_i$ 为余切丛纤维上的线性函数。再设 $W$ 是第 \ref{amo} 节意义下 $N$ 的复化。则 $Y=T^*W$ 是 $M$ 的复化。$Y$ 的自然全纯辛形式仍可写成 $\Omega=\sum_i\d p_i\d q^i$（此时 $q^i$ 与 $p_i$ 是 $W$ 上及 $T^*W\to W$ 纤维上的局部全纯坐标），它是 $\omega$ 的解析延拓。若 $Y$ 的 $A$-模型存在，就可用它量子化 $M$，并把结果与利用 $M=T^*N$ 的余切丛描述所作的几何量子化比较。

先讨论 $Y$ 的形变量子化，即讨论 $\Hom(\B_\cc,\B_\cc)$。余切丛有极大约束力的伸缩对称。为利用此对称，方便的是稍作推广，取 $\Omega=\sum_i\d p_i\d q^i/\hbar$（在得知可以令 $\hbar=1$ 后即可令之）。有一个伸缩对称使 $q^i$ 不变，而 $p_j$ 与 $\hbar$ 的次数为 1。

这个伸缩对称是 $\C^*$ 在这些变量上的作用。我们将假定形变量子化应以 $\C^*$-不变的方式施行。若 $\C^*$ 的某个 $\UU(1)$ 子群是底层 $\sigma$-模型的对称（所关心的 $A$-模型由该 $\sigma$-模型导出），则确乎如此。（在将别处讨论的几何 Langlands 应用中这是对的 \cite{GaW}。）全纯性会把这样的 $\UU(1)$ 对称提升为 $\C^*$。若没有底层 $\UU(1)$ 对称，经由 $\B_\cc$ 的形变量子化可能不是 $\C^*$-不变的。

例如设 $W=\C^n$、$Y=\C^{2n}$，使坐标 $q^i,p_j$ 可整体定义。Poisson 括号为 $\{p_i,q^j\}=\delta_i^j$、$\{q^i,q^j\}=\{p_i,p_j\}=0$。要求量子换位子与 Poisson 括号在 $\hbar^2$ 阶修正以内一致，并施加伸缩对称，则量子换位子最一般的形式是
\begin{align}\label{genform}
[\h p_i,\h q^j]&=\i\hbar\delta^j_i\cr
[\h p_i,\h p_j]&=\hbar^2f_{ij}(\h q^k)\cr
[\h q^i,\h q^j]&=0.\end{align}
Jacobi 恒等式 $[\h p_i,[\h p_j,\h p_k]]+\cdots=0$ 蕴含 $\partial_if_{jk}+\partial_jf_{ki}+\partial_kf_{ij}=0$，故存在函数 $g_k$ 使 $f_{ij}=\partial_ig_j-\partial_jg_i$。于是把 $\h p_i$ 换成 $\h p_i'=\h p_i-\i\hbar g_i(q^k)$（此代换与伸缩关系相容），可把式 \eqref{genform} 化为
\begin{align}\label{genformtwo}
[\h p_i',\h q^j]&=\i\hbar\delta^j_i\cr
[\h p_i',\h p_j']&=0\cr
[\h q^i,\h q^j]&=0.\end{align}
因此，$p_i$ 与 $q^j$ 的多项式函数（如第 \ref{abranes} 节所述，这是我们唯一允许的函数）的量子形变代数，就是把 $q^j$ 表示为乘算符、$p_i$ 表示为 $-\i\hbar\frac{\partial}{\partial q^i}$ 所得的代数。当然，最后这个公式在初等量子力学中很熟悉。有了对 $\h q^i$ 与 $\h p_j$ 的这一刻画，结合代数 $\Hom(\B_\cc,\B_\cc)$ 可刻画为 $\C^n$ 上多项式全纯微分算符的代数。在建立这个与微分算符的关系时，重要的是只允许对 $p$ 的依赖为多项式的函数；下文将始终施加此条件。类似地也只允许对 $q$ 的多项式依赖。于是 $\h p$ 与 $\h q$ 的函数对应多项式全纯微分算符。

现在考虑更一般的 $W$。局部上可化归前一情形：用 $W$ 的开集 $\U\subset W$（选得使 $\U$ 同构于 $\C^n$ 中一个区域）代替 $W$，用 $T^*\U$ 代替 $Y$。把靶空间 $Y=T^*W$ 换成开集（此处为 $T^*\U$），在 $\sigma$-模型中一般不是好操作。不过，由于 $\sigma$-模型微扰论涉及靶空间中常值映射附近的小涨落，限制到开集在 $\sigma$-模型微扰论中没有问题。通常这只能给出微扰信息，但我们处在不寻常的情形：这已经足够，因为伸缩对称意味着到 $\hbar^2$ 阶的微扰论给出全部答案。我们选择形如 $T^*\U$ 的开集，是因为这样的集合在伸缩下不变，从而伸缩论证适用。因此可以把 $\Hom(\B_\cc,\B_\cc)$ 视为 $W$ 上的层：对任何开集 $\U\subset W$，联系以 $T^*\U$ 上的 $\Hom(\B_\cc,\B_\cc)$。应当承认，尽管这一步在数学上说得通且显得很自然，它在物理上是否完全自然并不清楚。最好能为这一步找到更自然的物理解释，或在后续论证中设法避开它。

若取 $\U$ 小到同构于 $\C^n$ 中的开球，事情就简单了。一旦限制到这样的 $\U$ 对应的 $T^*\U$，坐标 $q_i,p^j$ 可整体定义，并可像前面那样提升为表示 $(\B_\cc,\B_\cc)$ 弦的算符。故上述分析适用于任何这样的集合。

现在用小开集 $\U_\alpha$ 覆盖 $W$，每个同构于 $\C^n$ 中的开集，且交集拓扑平凡。于是 $Y$ 被开集 $T^*\U_\alpha$ 覆盖。对每个 $\alpha$，取 $\U_\alpha$ 的局部坐标 $q^i_\alpha$ 及余切丛的纤维坐标 $p_{j,\alpha}$。在交集 $T^*(\U_\alpha\cap\U_\beta)$ 中，坐标之间的关系为
\begin{align}\label{relco}
q^i_\alpha&=h^i_{\alpha\beta}(q^k_\beta)\cr
p_{i,\alpha}&=M_{\alpha\beta}(q)^k_i\,p_{k,\beta},\end{align}
其中 $h^i_{\alpha\beta}$ 是某些全纯函数，$M_{\alpha\beta}{}^i_j$ 是 $\partial_ih^j_{\alpha\beta}$ 的逆矩阵。

量子层面，伸缩对称不允许 $q$ 的转移公式有任何改变，故必须有 $\h q^i_\beta=h^i_{\alpha\beta}(\h q^k)$，其中 $q$ 的函数与之前相同。但对 $p$，伸缩对称允许 $\hbar$ 阶修正，于是公式变为 $\h p_{i,\alpha}=M_{\alpha\beta}(\h q)^k_i\h p_{k,\beta}+\i\hbar r_{i,\alpha\beta}(\h q)$，其中 $r_{i,\alpha\beta}$ 为某些全纯函数。然而为满足 $[\h p_{i,\alpha},\h p_{j,\beta}]=0$，需要 $\partial_ir_{j,\alpha\beta}-\partial_jr_{i,\alpha\beta}=0$，因此（由于 $\U_\alpha\cap\U_\beta$ 拓扑平凡）存在全纯函数 $w_{\alpha\beta}$ 使 $r_{j,\alpha\beta}=\partial_jw_{\alpha\beta}(\h q)$。于是 $\h p$ 的变换为
\be\label{transp}\h p_{i,\alpha}=M_{\alpha\beta}(\h q)^k_i\h p_{k,\beta}+\i\hbar\partial_iw_{\alpha\beta}(\h q).\ee
若在区域 $\U_\alpha\cap\U_\beta$ 中只用一组坐标 $q^i$（而不用两组坐标 $q^i_\alpha,q^i_\beta$），此公式会简单得多。此时矩阵 $M$ 换成 1，公式简化为
\be\label{transpo}\h p_{i,\alpha}=\h p_{i,\beta}+\i\hbar\partial_iw_{\alpha\beta}.\ee
这就分离出了 $\h p$ 变换中的非经典部分。

立即可得 $w_{\alpha\beta}=-w_{\beta\alpha}$；三重交上的相容性告诉我们 $\partial_k(w_{\alpha\beta}+w_{\beta\gamma}+w_{\gamma\alpha})=0$，故 $\phi_{\alpha\beta\gamma}=w_{\alpha\beta}+w_{\beta\gamma}+w_{\gamma\alpha}$ 是复常数。在四重交上，这些常数自动满足
\be\label{cocrel}\phi_{\alpha\beta\gamma}-\phi_{\beta\gamma\delta}+\phi_{\gamma\delta\alpha}-\phi_{\delta\alpha\beta}=0,\ee
因而定义出 $H^2(W,\C)$ 的一个元素 $[\phi]$。

若 $[\phi]$ 实际取值于 $2\pi\i\cdot H^2(W,\Z)$，即（在可能把 $w_{\alpha\beta}$ 平移常数之后）每个 $\phi_{\alpha\beta\gamma}$ 都是 $2\pi\i$ 的整数倍，则量 $S_{\alpha\beta}=\exp(w_{\alpha\beta})$ 满足 $S_{\alpha\beta}S_{\beta\gamma}S_{\gamma\alpha}=1$，可视为复线丛 $\L\to W$ 的转移函数。具体地说，$\L$ 如下：在每个 $\U_\alpha$ 中，$\L$ 的截面就是函数 $\psi_\alpha$；但不同开集 $\U_\alpha$ 与 $\U_\beta$ 中的函数由 $\psi_\alpha=S_{\alpha\beta}\psi_\beta$ 相联系。在此情形，上述公式可解释为：$q$ 与 $p$ 的函数的量子化给出作用在 $\L$ 截面上的全纯微分算符。换言之，设 $D_\alpha$ 与 $D_\beta$ 分别是作用在相应开集 $\U_\alpha$ 与 $\U_\beta$ 中函数上的全纯微分算符。为把它们解释为 $W$ 上整体定义的全纯微分算符限制到这些区域，需要 $D_\alpha(S_{\alpha\beta}\psi)=S_{\alpha\beta}D_\beta\psi$，即 $D_\alpha=S_{\alpha\beta}D_\beta S^{-1}_{\alpha\beta}$。特别地，在区域 $\beta$ 中为 $\h p_i=-\i\hbar\frac{\partial}{\partial q^i}$ 的算符，在区域 $\alpha$ 中将是 $\h p_i+\i\hbar\partial_i\log S_{\alpha\beta}=\h p_i+\i\hbar\partial_iw_{\alpha\beta}$，与式 \eqref{transp} 相符。

这有如下推广。\footnote{\label{zany}我们将描述的情形出现在量子化的重要例子中，但并不完全一般。例如，若 $W$ 紧，且 $[\phi]$ 的 Hodge 分解含有 $(2,0)$ 型部分，则下文所述不适用。}设 $\L\to W$ 是全纯线丛，并设在各开集 $\U_\alpha$ 中取定 $\L$ 的平凡化后，$\L$ 可由 $\U_\alpha\cap\U_\beta$ 中的转移函数 $e^{x_{\alpha\beta}}$ 描述。设 $w_{\alpha\beta}=\lambda x_{\alpha\beta}$，$\lambda$ 为某复数。若 $\lambda$ 为整数，则函数 $e^{w_{\alpha\beta}}=e^{\lambda x_{\alpha\beta}}$ 是某个全纯线丛的转移函数，即 $\L^\lambda$。此时上一段的讨论适用，$q$ 与 $p$ 的函数的量子化是作用在 $\L^\lambda$ 截面上的全纯微分算符。对一般的复 $\lambda$，线丛 $\L^\lambda$ 并不存在，但这只是因为三重交 $\U_\alpha\cap\U_\beta\cap\U_\gamma$ 中转移函数之积 $e^{\lambda(x_{\alpha\beta}+x_{\beta\gamma}+x_{\gamma\alpha})}$ 是 $c$-数而不等于 1。由于微分算符与 $c$-数交换，无论 $\L^\lambda$ 是否作为线丛存在，作用在 $\L^\lambda$ 上的微分算符这个概念都有意义。于是 $q$ 与 $p$ 的量子化函数自然地理解为作用在 $\L^\lambda$ 上的微分算符。

更一般地，设 $\L_i\to W$ 是全纯线丛，由全纯函数 $x_{i,\alpha\beta}$ 给出的转移函数 $e^{x_{i,\alpha\beta}}$ 描述。若 $w_{\alpha\beta}=\sum_i\lambda_ix_{i,\alpha\beta}$，其中 $\lambda_i$ 为某些整数，则 $\otimes_i\L_i^{\lambda_i}$ 作为线丛有意义，$p$ 与 $q$ 的量子化函数是作用在 $\otimes_i\L_i^{\lambda_i}$ 上的微分算符。若 $w_{\alpha\beta}=\sum_i\lambda_ix_{i,\alpha\beta}$ 而 $\lambda_i$ 为更一般的复系数，则 $\otimes_i\L_i^{\lambda_i}$ 作为线丛没有意义，但作用在 $\otimes_i\L_i^{\lambda_i}$ 上的微分算符仍有意义，且 $p$ 与 $q$ 的量子化函数自然地是作用在 $\otimes_i\L_i^{\lambda_i}$ 上的全纯微分算符。

我们将松散地把 $\otimes_i\L_i^{\lambda_i}$ 称为线丛的复幂（尽管更准确地说，是这种复幂的张量积）。我们所描述的情形对量子化很重要，但如脚注 \ref{zany} 所述，它不是最一般的可能性。

一个重要细节是：若 $\RR$ 是容许平坦联络的全纯线丛，则作用在 $\otimes_i\L_i^{\lambda_i}$ 上的全纯微分算符与作用在 $\RR\otimes_i\L_i^{\lambda_i}$ 上的全纯微分算符相同。这是因为平坦线丛可由其 $c$-数整体和乐刻画；而微分算符与 $c$-数交换，故全纯微分算符对 $\otimes_i\L_i^{\lambda_i}$ 与 $\RR\otimes_i\L_i^{\lambda_i}$ 之别不敏感。

结果表明，在 $Y=T^*W$ 的情形，$\Hom(\B_\cc,\B_\cc)$ 就是作用在线丛复幂上的微分算符层。证明这一点的一个快捷办法是方便地选取 $A$-膜 $\B_0$，并利用 $\Hom(\B_\cc,\B_\cc)$ 可作用在 $\Hom(\B_0,\B_\cc)$ 上这一事实。为此我们将利用 $(B,A,A)$ 型拉格朗日膜。

\begin{yizhu}[线丛的“复幂”是什么意思]
$\L=\otimes_i\L_i^{\lambda_i}$ 在 $\lambda_i$ 非整数时并非真正的线丛：局部转移函数 $e^{\lambda_i x_{\alpha\beta}}$ 仍是 $U(1)$-值的，但三重交不再自动满足上闭链条件，故整体可能“粘不起来”。原文仍使用它，是因为全纯微分算符与这些 $c$-数因子交换，所以“作用在 $\L$ 的截面上”这句话在局部仍有确切含义；在 $W$ 紧且 $[\phi]$ 含 $(2,0)$ 部分时，原文脚注已提醒这一描述会失效。[来源：I 级]
\end{yizhu}
